\section*{Reading Paths}
\addcontentsline{toc}{section}{Reading Paths}
\label{sec:front-matter-reading-paths}

The manual is written for several entry points:

\begin{itemize}
\item first contact;
\item GUI modeling;
\item terminal execution;
\item component development;
\item architecture;
\item tests;
\item scientific simulation;
\item biological extensions.
\end{itemize}

Figure~\ref{fig:manual-reading-map} shows the recommended progression through the book.

% FIGURE-SPEC-BEGIN
% ID: fig:manual-reading-map
% Source asset: figs/generated/manual-reading-map
% Caption: Recommended progression through the manual from first contact to advanced topics.
% Accessibility: A structured map that highlights the reading order for user, developer, and advanced technical paths.
% FIGURE-SPEC-END
\begin{figure}[htbp]
\centering
\figureplaceholder{figs/generated/manual-reading-map}
\caption{Recommended progression through the manual from first contact to advanced topics.}
\label{fig:manual-reading-map}
\end{figure}

Figure~\ref{fig:manual-audience-matrix} maps audiences to the chapters that matter most to them.

% FIGURE-SPEC-BEGIN
% ID: fig:manual-audience-matrix
% Source asset: figs/generated/manual-audience-matrix
% Caption: Audience-to-chapter matrix for the current manual structure.
% Accessibility: A matrix that helps the reader identify the chapters most relevant to a given role or goal.
% FIGURE-SPEC-END
\begin{figure}[htbp]
\centering
\figureplaceholder{figs/generated/manual-audience-matrix}
\caption{Audience-to-chapter matrix for the current manual structure.}
\label{fig:manual-audience-matrix}
\end{figure}
